\cotefile{157-3}
\scanpage{237}
\mspage{187}

\[
\operatorname{Hom}_{\pi}(\pi_{n-1}(Z,x),\underset{\substack{\Vert\\\inserted{F_{\unsure{x|n}}}}}{\deleted{\unsure{A}}M})=0
\]

et ceci \uline{\unsure{quel|quels} que \unsure{soit|soient} \unsure{le|les}} $\pi$-\unsure{module|modules} $M$, puisque $\pi=\pi_1(X,x)\simeq\pi_1(Y,y)$. Cela implique (en prenant p.~ex. $M=\pi_{n-1}(Z,x)$ !) que $\pi_{n-1}(Z,x)=0$. (\deleted{\unsure{Il}}Mais en faisant attention dans le cas $n=2$, qu'on sait déjà que $\pi_1(Z,\unsure{x})$ est abélien).

Donc le th. rectifié \unsure{page} $178$ a bien l'air vrai. Il est un peu triste, puisque dans chacune des implications

\[
a_n\Longrightarrow b_n\Longrightarrow c_n\Longrightarrow a_{n-1}
\]

on perd un chouïa, et en fin de compte on sort avec les magnifiques implications

\[
a_n\Longrightarrow a_{n-1}.
\]

De toutes façons, il est $\pm$ clair que le ``bon'' localiseur $W_n$, celui qui doit définir les ``$n$-types d'homotopie'' comme objets de $\operatorname{Hot}_{W_n}$, est défini par la condition $\unsure{a_n}$.

2. \uline{Critère de $W_\infty$-asphéricité}

Je vais essayer de faire mieux quand $Y$ est réduit à un point. Cette fois, je
% archive sheet number: 237
